\section{Complex Scalar Field and Global Phase Symmetry}
\label{sec:complex_scalar_phase_symmetry}

A complex scalar field provides the standard example of an internal continuous symmetry. The symmetry rotates the phase of the field by the same angle at every spacetime point. Noether's theorem then produces a conserved current that is independent of spacetime translation symmetry.

\subsection{Complex scalar field}

Let

\[
\phi:
\mathbb{R}^{1,3}
\longrightarrow
\mathbb{C}.
\]

The complex conjugate \(\phi^{*}\) is treated as an independent variable during variation. Consider the free density

\[
\boxed{
\mathcal{L}
=
\partial_{\mu}\phi^{*}
\partial^{\mu}\phi
-
\mu^{2}\phi^{*}\phi
}.
\]

The action is real because

\[
\left(
\partial_{\mu}\phi^{*}
\partial^{\mu}\phi
\right)^{*}
=
\partial_{\mu}\phi^{*}
\partial^{\mu}\phi,
\]

and

\[
\left(
\phi^{*}\phi
\right)^{*}
=
\phi^{*}\phi.
\]

\subsection{Field equations}

Variation with respect to \(\phi^{*}\) gives

\[
\frac{\partial\mathcal{L}}
{\partial\phi^{*}}
=
-\mu^{2}\phi,
\]

and

\[
\frac{\partial\mathcal{L}}
{\partial(\partial_{\mu}\phi^{*})}
=
\partial^{\mu}\phi.
\]

Therefore

\[
\boxed{
\Box\phi+\mu^{2}\phi=0
}.
\]

Variation with respect to \(\phi\) gives the conjugate equation

\[
\boxed{
\Box\phi^{*}+\mu^{2}\phi^{*}=0
}.
\]

\subsection{Global phase symmetry}

Consider

\[
\boxed{
\phi
\longrightarrow
\phi'
=
e^{-i\alpha}\phi
},
\]

\[
\boxed{
\phi^{*}
\longrightarrow
\phi'^{*}
=
e^{i\alpha}\phi^{*}
},
\]

where \(\alpha\) is constant.

The potential term is invariant:

\[
\phi'^{*}\phi'
=
e^{i\alpha}
e^{-i\alpha}
\phi^{*}\phi
=
\phi^{*}\phi.
\]

Because \(\alpha\) is constant,

\[
\partial_{\mu}\phi'
=
e^{-i\alpha}\partial_{\mu}\phi,
\]

and

\[
\partial_{\mu}\phi'^{*}
=
e^{i\alpha}\partial_{\mu}\phi^{*}.
\]

Thus

\[
\partial_{\mu}\phi'^{*}
\partial^{\mu}\phi'
=
\partial_{\mu}\phi^{*}
\partial^{\mu}\phi.
\]

Therefore

\[
\boxed{
\mathcal{L}'=\mathcal{L}
}.
\]

\subsection{Infinitesimal form}

For small \(\alpha=\varepsilon\),

\[
e^{-i\varepsilon}
=
1-i\varepsilon+O(\varepsilon^{2}),
\]

so

\[
\boxed{
\delta\phi
=
-i\varepsilon\phi
},
\qquad
\boxed{
\delta\phi^{*}
=
i\varepsilon\phi^{*}
}.
\]

The coordinate variation vanishes:

\[
\xi^{\mu}=0.
\]

The field generators are

\[
\Psi_{\phi}
=
-i\phi,
\qquad
\Psi_{\phi^{*}}
=
i\phi^{*}.
\]

\subsection{Noether current}

The derivative momenta are

\[
\frac{\partial\mathcal{L}}
{\partial(\partial_{\mu}\phi)}
=
\partial^{\mu}\phi^{*},
\]

and

\[
\frac{\partial\mathcal{L}}
{\partial(\partial_{\mu}\phi^{*})}
=
\partial^{\mu}\phi.
\]

Noether's formula gives

\[
\begin{aligned}
j^{\mu}
&=
\frac{\partial\mathcal{L}}
{\partial(\partial_{\mu}\phi)}
\left(
-i\phi
\right)
+
\frac{\partial\mathcal{L}}
{\partial(\partial_{\mu}\phi^{*})}
\left(
i\phi^{*}
\right)
\\
&=
-i\phi\partial^{\mu}\phi^{*}
+
i\phi^{*}\partial^{\mu}\phi.
\end{aligned}
\]

Hence

\[
\boxed{
j^{\mu}
=
i
\left(
\phi^{*}\partial^{\mu}\phi
-
\phi\partial^{\mu}\phi^{*}
\right)
}.
\]

\subsection{Verification of conservation}

Differentiate:

\[
\begin{aligned}
\partial_{\mu}j^{\mu}
&=
i
\partial_{\mu}
\left(
\phi^{*}\partial^{\mu}\phi
-
\phi\partial^{\mu}\phi^{*}
\right)
\\
&=
i
\left[
(\partial_{\mu}\phi^{*})
(\partial^{\mu}\phi)
+
\phi^{*}\Box\phi
-
(\partial_{\mu}\phi)
(\partial^{\mu}\phi^{*})
-
\phi\Box\phi^{*}
\right].
\end{aligned}
\]

The gradient products cancel:

\[
\partial_{\mu}j^{\mu}
=
i
\left(
\phi^{*}\Box\phi
-
\phi\Box\phi^{*}
\right).
\]

Using

\[
\Box\phi
=
-\mu^{2}\phi,
\]

and

\[
\Box\phi^{*}
=
-\mu^{2}\phi^{*},
\]

we obtain

\[
\partial_{\mu}j^{\mu}
=
i
\left(
-\mu^{2}\phi^{*}\phi
+
\mu^{2}\phi\phi^{*}
\right)
=
0.
\]

Therefore

\[
\boxed{
\partial_{\mu}j^{\mu}=0
}.
\]

\subsection{Charge density and charge}

With

\[
j^{\mu}
=
(c\rho,\mathbf{j}),
\]

the charge density is

\[
\boxed{
\rho
=
\frac{i}{c}
\left(
\phi^{*}\partial^{0}\phi
-
\phi\partial^{0}\phi^{*}
\right)
}.
\]

Since

\[
\partial^{0}
=
\frac{1}{c}
\frac{\partial}{\partial t},
\]

this becomes

\[
\boxed{
\rho
=
\frac{i}{c^{2}}
\left(
\phi^{*}\phi_t
-
\phi\phi_t^{*}
\right)
}.
\]

The conserved charge is

\[
\boxed{
Q
=
\int
\rho
\,\mathrm{d}^{3}x
}.
\]

The overall normalization of \(j^{\mu}\) can be multiplied by a constant together with a corresponding normalization of the symmetry generator. The conservation law is unaffected.

\subsection{Plane-wave example}

Let

\[
\phi(x)
=
A
e^{-ik_{\mu}x^{\mu}},
\]

where \(A\) is constant. Then

\[
\partial^{\mu}\phi
=
-ik^{\mu}\phi,
\]

and

\[
\partial^{\mu}\phi^{*}
=
ik^{\mu}\phi^{*}.
\]

The current is

\[
\begin{aligned}
j^{\mu}
&=
i
\left[
\phi^{*}
\left(
-ik^{\mu}\phi
\right)
-
\phi
\left(
ik^{\mu}\phi^{*}
\right)
\right]
\\
&=
2k^{\mu}|A|^{2}.
\end{aligned}
\]

Thus

\[
\boxed{
j^{\mu}
=
2k^{\mu}|A|^{2}
}.
\]

The current is parallel to the wave four-vector.

\subsection{Two-real-field interpretation}

Write

\[
\phi
=
\frac{1}{\sqrt{2}}
\left(
\phi_1+i\phi_2
\right).
\]

Then

\[
\mathcal{L}
=
\frac{1}{2}
\partial_{\mu}\phi_1
\partial^{\mu}\phi_1
+
\frac{1}{2}
\partial_{\mu}\phi_2
\partial^{\mu}\phi_2
-
\frac{1}{2}
\mu^{2}
\left(
\phi_1^{2}+\phi_2^{2}
\right).
\]

The phase transformation is a rotation:

\[
\delta\phi_1
=
\varepsilon\phi_2,
\qquad
\delta\phi_2
=
-\varepsilon\phi_1,
\]

up to the sign convention for the parameter. The current becomes

\[
\boxed{
j^{\mu}
=
\phi_2\partial^{\mu}\phi_1
-
\phi_1\partial^{\mu}\phi_2
}.
\]

Thus the conserved charge measures rotational motion in the internal \((\phi_1,\phi_2)\) plane.

\subsection{Interactions that preserve the symmetry}

The most general potential depending only on

\[
|\phi|^{2}
=
\phi^{*}\phi
\]

preserves the global phase symmetry:

\[
\mathcal{L}
=
\partial_{\mu}\phi^{*}
\partial^{\mu}\phi
-
V(\phi^{*}\phi).
\]

The equations are nonlinear when \(V\) is nonlinear, but the same current remains conserved because the potential is phase invariant.

For example,

\[
V
=
\mu^{2}\phi^{*}\phi
+
\lambda
\left(
\phi^{*}\phi
\right)^{2}
\]

preserves the symmetry.

A term such as

\[
\kappa
\left(
\phi^{2}+\phi^{*2}
\right)
\]

does not remain invariant under a general phase rotation and explicitly breaks the \(U(1)\) symmetry.

\subsection{Why a local phase requires a gauge field}

Let

\[
\alpha=\alpha(x).
\]

Then

\[
\partial_{\mu}
\left(
e^{-i\alpha(x)}\phi
\right)
=
e^{-i\alpha(x)}
\left[
\partial_{\mu}\phi
-
i
(\partial_{\mu}\alpha)\phi
\right].
\]

The extra derivative of the parameter prevents invariance of the kinetic term. A gauge field and a covariant derivative are introduced precisely to cancel this additional term. This observation prepares the transition from global internal symmetry to electromagnetic gauge freedom.

\subsection{Common mistakes}

Typical errors include:

\begin{enumerate}
    \item varying \(\phi\) but not \(\phi^{*}\);
    \item assigning the same phase to \(\phi\) and \(\phi^{*}\);
    \item forgetting that \(\alpha\) must be constant for the ordinary derivative term to remain invariant;
    \item losing the factor of \(i\) in the current;
    \item concluding that a complex field contains only one real degree of freedom;
    \item assuming that every interaction breaks the phase symmetry rather than checking whether it depends only on \(|\phi|^{2}\).
\end{enumerate}

\subsection{What has been established}

The complex scalar action is invariant under the global phase transformation

\[
\phi
\longrightarrow
e^{-i\alpha}\phi.
\]

Noether's theorem produces the conserved current

\[
j^{\mu}
=
i
\left(
\phi^{*}\partial^{\mu}\phi
-
\phi\partial^{\mu}\phi^{*}
\right).
\]

The current is conserved on solutions of the complex Klein--Gordon equation. In terms of two real fields, the symmetry is an ordinary rotation in internal field space. Promoting the phase parameter to \(\alpha(x)\) fails with ordinary derivatives and motivates the gauge-covariant structures introduced in the next chapter.
